\subsection*{Задание 1:}
\lstinputlisting[caption={Исходный код программы (C)}]{1.c}
\noindent \indent
Данное задание можно решить просто пройдясь по всем элементам массива и сравнить текущий индекс элемента
с индексом котрый возвращает цепочка из трёх элементов.
\newpage
\subsection*{Задание 2:}
% \begin{figure}[H]
% 	\centering
% 	\includegraphics[width=\textwidth, height=0.9\textheight, keepaspectratio]{2.pdf}
% 	\caption{Схема алгоритма задачи}
% \end{figure}
% \newpage
\lstinputlisting[caption={Исходный код программы (C)}]{2.c}
\noindent \indent
Это задание можно решить сразу при вводе, проверя прямой ход пешки или наискосок если там стоит вражеская пешка.
\newpage
\subsection*{Задание 3:}
% \begin{figure}[H]
% 	\centering
% 	\includegraphics[width=\textwidth, height=0.9\textheight, keepaspectratio]{3.pdf}
% 	\caption{Схема алгоритма задачи}
% \end{figure}
% \newpage
\lstinputlisting[caption={Исходный код программы (C)}]{3.c}
\noindent \indent
Для решения этого задания надо вывести лабиринт в котором k точек старта приводят к выходу, а остальные нет.
Для этого можно завести счётчик и сводить к выходу влево-вверх пока счётчик меньше k, а иначе сводидть вправо-вниз
и развернуть последнюю влево.
% \newpage
% \subsection*{Задание 4:}
% \lstinputlisting[caption={Исходный код программы (C)}]{4.c}
% \noindent \indent
% Суть данного задания - вычислить размер фотографии $w \cdot h$, где $w$ - сумма $w_i$, а $h$ - максимальное $h_i$
% для каждого друга, так как один из друзей фотографирует остальных. Для решения задачи суммируем w, находим максимальный
% и следующий после него $h$ и вычисляем размер, вычитая из суммы $w_i$ и умножаешь на максимальное (или на следующее максимальное).
% \newpage
% \subsection*{Задание 5:}
% \begin{figure}[H]
% 	\centering
% 	\includegraphics[width=\textwidth, height=0.9\textheight, keepaspectratio]{5.pdf}
% 	\caption{Схема алгоритма задачи}
% \end{figure}
% \newpage
% \lstinputlisting[caption={Исходный код программы (Python)}]{5.py}
% \noindent \indent
% Суть данного задания - найти кол-во пар элементов, которые можно удалить из массива чтобы среднее арифметическое не поменялось.
% Для реализации алгоритма находим целевую сумму чисел, если сумма вещественная - то из входных данных мы не найдём подходящей пары.
% В цикле проходим по элементам и находим элемент, который нам нужен $need = target - x$, добавляем к счётчику число из словаря $ans += count[need]$,
% и увеличиваем значение текущего эл. в словаре $count[x] += 1$
% \newpage